\subsection{Deductive Logic Agent và deduction loop}
\label{sec:deductive_agent}
\begin{sloppypar}

Mục này mô tả \texttt{DeductiveAgent.solve}, vòng lặp đưa tác nhân từ trạng
thái khởi đầu tới khi ván kết thúc. Mỗi vòng lặp là một lượt của đường ống ở
Mục~\ref{sec:agent_pipeline}: tìm một phán quyết chứng minh được, nộp nó qua
Game Engine, nhận lời khai mới và lặp lại trên $\mathit{KB}$ lớn hơn.

\paragraph{Vòng lặp suy luận.}
Thuật toán~\ref{alg:deduction_loop} trình bày vòng lặp theo đúng cài đặt.
Hàm con \textsc{NextMove} là \texttt{next\_move}: kiểm tra $\mathit{KB}_t$ nhất
quán, rồi duyệt các ô chưa chứng minh theo thứ tự định danh và trả về phán
quyết của ô đầu tiên có nhãn \texttt{CRIMINAL} hoặc \texttt{INNOCENT}
(Mục~\ref{sec:entailment_classification}).

\begin{algorithm}[htbp]
  \caption{Vòng lặp suy luận \texttt{DeductiveAgent.solve}}
  \label{alg:deduction_loop}
  \begin{algorithmic}[1]
    \Require Game Engine $E$ đã nạp bản đồ; bộ kiểm tra \textsc{HasSolution}
    \Ensure Trạng thái dừng, ảnh chụp công khai cuối cùng và trace
    \State $\mathit{trace} \gets [\,]$;\quad $S \gets E.\textsc{GetPublicState}()$
    \Loop
      \State $(\mathit{status}, v, Q) \gets \textsc{NextMove}(S)$
        \Comment{$v = \langle i, s\rangle$ hoặc rỗng; $Q$: các truy vấn đã hỏi}
      \If{$v$ rỗng}
        \If{$\mathit{status} = \texttt{UNKNOWN}$}
          \State \Return $(\texttt{STUCK}, S, \mathit{trace})$
        \EndIf
        \State \Return $(\mathit{status}, S, \mathit{trace})$
          \Comment{\texttt{SOLVED} hoặc \texttt{INCONSISTENT}}
      \EndIf
      \State $r \gets E.\textsc{SubmitVerdict}(i, s)$
      \If{$r.\mathit{outcome} \ne \texttt{ACCEPTED}$}
        \State \textbf{raise} \texttt{MoveRejected}$(v, r.\mathit{outcome})$
      \EndIf
      \State $A \gets$ định danh các lời khai trong $S$
      \State $N \gets$ định danh các lời khai trong $r.\mathit{state}$ không thuộc $A$
      \State thêm $\textsc{Step}(|\mathit{trace}|+1,\ v,\ Q,\ A,\ N)$ vào $\mathit{trace}$
      \State $S \gets r.\mathit{state}$
    \EndLoop
  \end{algorithmic}
\end{algorithm}

\paragraph{Gửi verdict và cập nhật $\mathit{KB}$.}
Tác nhân không tự sửa trạng thái ván chơi. Nó chỉ gọi
\texttt{submit\_verdict(i, s)}, và Game Engine tự kiểm tra lại phán quyết bằng
hai truy vấn của mình (Mục~\ref{sec:verdict_acceptance}). Khi kết quả là
\texttt{ACCEPTED}, engine ghi $\langle i, s\rangle$ vào
\texttt{proved\_statuses} và lật mọi lời khai do ô $i$ sở hữu theo
\eqref{eq:reveal}, rồi trả về ảnh chụp công khai mới. Tác nhân dùng ảnh chụp
này làm đầu vào cho vòng kế tiếp. Vì \texttt{DPLLChecker} dựng lại CNF từ
ảnh chụp ở mỗi lần gọi, việc ``cập nhật $\mathit{KB}$'' không cần thao tác
riêng: $\mathit{KB}_{t+1}$ tự động gồm mệnh đề đơn vị của phán quyết vừa được
chấp nhận và các mệnh đề của lời khai vừa lật. Theo Mục~\ref{sec:reveal_protocol},
$\mathit{KB}_{t+1}$ chỉ thêm mệnh đề so với $\mathit{KB}_t$, nên ô nào đã suy ra
được ở vòng trước vẫn suy ra được ở vòng sau.

Nếu engine từ chối, tác nhân ném ngoại lệ \texttt{MoveRejected} kèm phán quyết
và kết quả. Trường hợp này chỉ xảy ra khi tác nhân và engine dùng hai bộ kiểm
tra cho kết quả khác nhau, tức là có lỗi cài đặt. Vì vậy tác nhân dừng hẳn
thay vì bỏ qua và thử ô khác; API chuyển ngoại lệ này thành mã lỗi HTTP $409$.

\paragraph{Ba trạng thái dừng.}
Vòng lặp kết thúc ở một trong ba trạng thái của kiểu \texttt{SolveStatus}:
\begin{itemize}
  \item \texttt{SOLVED}: mọi ô đã được chứng minh. \textsc{NextMove} xác nhận
    $\mathit{KB}_t$ nhất quán, thấy $U_t$ rỗng và trả trạng thái này. Vì mỗi
    phán quyết được chấp nhận đều là hệ quả của $\mathit{KB}_t$, cấu hình cuối
    cùng chính là $\sigma^\star$ dù tác nhân không bao giờ đọc $\sigma^\star$.
  \item \texttt{STUCK}: $\mathit{KB}_t$ nhất quán, còn ô chưa chứng minh nhưng
    mọi ô đều có nhãn \texttt{UNKNOWN}. Tác nhân dừng thay vì đoán, đúng với
    luật không đoán. Trạng thái này cho biết bản đồ vi phạm điều kiện mà
    \texttt{\_check\_solvable\_without\_guessing} kiểm tra
    (Mục~\ref{sec:uniqueness_checking}), nên không xảy ra trên sáu bản đồ đã
    được kiểm chứng.
  \item \texttt{INCONSISTENT}: $\mathit{KB}_t$ không thỏa được. Bản đồ đã kiểm
    chứng có $\sigma^\star$ thỏa mọi lời khai nên cũng không rơi vào trạng thái
    này; nó tồn tại để tác nhân báo lỗi rõ ràng khi được dùng trên một bản đồ
    hỏng. Ca kiểm thử
    \texttt{test\_classify\_all\_\allowbreak returns\_\allowbreak inconsistent\_\allowbreak on\_unsat\_kb} dựng
    một $\mathit{KB}$ mâu thuẫn và xác nhận tác nhân trả về đúng trạng thái
    này.
\end{itemize}
Vòng lặp luôn dừng: mỗi vòng không dừng thêm đúng một ô vào $\mathcal{P}_t$,
nên có tối đa $|\mathcal{C}| - |\mathcal{P}_0|$ vòng nộp phán quyết, cộng một
lần gọi \textsc{NextMove} cuối cùng. Mỗi lần gọi \textsc{NextMove} tốn tối đa
$1 + 2|U_t|$ truy vấn, nên tổng số truy vấn của tác nhân là $O(|\mathcal{C}|^2)$
trong trường hợp xấu nhất. Trên thực tế con số này nhỏ hơn nhiều vì
\textsc{NextMove} dừng ở ô đầu tiên chứng minh được.

\paragraph{Vết suy luận.}
Mỗi phán quyết được chấp nhận sinh một bản ghi \texttt{Step} gồm số thứ tự,
phán quyết, danh sách \texttt{Query} (giả thiết và kết quả thỏa được) mà
\textsc{NextMove} đã hỏi, các lời khai đang mở khi suy luận
(\texttt{active\_clue\_ids}) và các lời khai vừa được lật
(\texttt{revealed\_clue\_ids}). Trace chỉ chứa thông tin công khai, nên mỗi
bước có thể được người đọc kiểm tra lại chỉ bằng các lời khai đã mở. Endpoint
\texttt{POST /api/game/auto-solve} trả về trace này, và giao diện dùng nó để
phát lại từng bước.

Bảng~\ref{tab:market_trace} là trace đầy đủ của tác nhân trên bản đồ Khu chợ
$3\times3$, sinh bằng \texttt{DPLLChecker}. Ván bắt đầu với An (ô~$0$, A1) vô
tội và lời khai của An đã mở. Vì mỗi nhân vật giữ đúng một lời khai, phán quyết
được chấp nhận ở mỗi bước luôn lật lời khai của chính nhân vật đó; cột
\emph{Mở} đếm số lời khai dùng được \emph{trước} khi phán quyết được nộp. Cột
\emph{Truy vấn} đếm cả truy vấn kiểm tra nhất quán, còn cột \emph{Lý do} diễn
giải bằng lời phép suy ra mà truy vấn bác bỏ xác nhận, gọi lời khai theo tên
người nói. Nội dung các lời khai được liệt kê trong
Phụ lục~\ref{sec:map_specifications}.

\begin{table}[htbp]
  \centering\small
  \caption{Trace suy luận của \texttt{DeductiveAgent} trên bản đồ Khu chợ}
  \label{tab:market_trace}
  \begin{tabularx}{\textwidth}{clcc>{\raggedright\arraybackslash}X}
    \toprule
    \textbf{Bước} & \textbf{Phán quyết} & \textbf{Truy vấn} & \textbf{Mở}
      & \textbf{Lý do} \\
    \midrule
    1 & Chi (C1) vô tội     & 5 & 1
      & Bình (ô~1) còn \texttt{UNKNOWN}; An nói Chi vô tội \\
    2 & Dũng (A2) vô tội    & 5 & 2
      & Bình còn \texttt{UNKNOWN}; An nói Dũng vô tội \\
    3 & Bình (B1) tội phạm  & 3 & 3
      & Dũng: hàng 1 có đúng một tội phạm, mà An và Chi vô tội \\
    4 & Giang (B2) tội phạm & 3 & 4
      & Dũng: Giang cùng trạng thái với Bình \\
    5 & Hà (C2) vô tội      & 3 & 5
      & Chi nói Hà vô tội \\
    6 & Hùng (A3) tội phạm  & 3 & 6
      & Chi: Lan khác Hà nên Lan là tội phạm; Giang: bốn góc có số tội phạm
        chẵn, A1 và C1 vô tội, nên Hùng cũng là tội phạm \\
    7 & Khoa (B3) vô tội    & 3 & 7
      & Bình nói Khoa vô tội \\
    8 & Lan (C3) tội phạm   & 3 & 8
      & Chi: Lan khác Hà, mà Hà vô tội \\
    \bottomrule
  \end{tabularx}
\end{table}

Có thể đọc ra ba đặc điểm của vòng lặp từ bảng này. Thứ nhất, ở bước~1 và~2,
tác nhân phải hỏi về ô~$1$ trước (hai truy vấn đều thỏa được) rồi mới tới ô
chứng minh được, nên mỗi bước tốn $5$ truy vấn; từ bước~3 trở đi ô có định danh
nhỏ nhất luôn chứng minh được ngay, nên mỗi bước chỉ tốn $3$ truy vấn. Cộng
thêm truy vấn nhất quán cuối cùng xác nhận \texttt{SOLVED}, tác nhân gọi bộ
giải $5+5+6\cdot3+1 = 29$ lần, khớp với số liệu benchmark
(Bảng~\ref{tab:benchmark_raw}). Thứ hai, thứ tự định danh không trùng với thứ
tự sóng suy luận. Lan (ô~$8$) đã suy ra được từ bước~2, ngay sau khi lời khai của
Chi được lật, nhưng chỉ được nộp ở bước cuối vì có định danh lớn nhất; ở bước~6,
suy luận về Hùng phải đi qua chính trạng thái của Lan dù Lan chưa được nộp.
Thứ ba, Bình là ô duy nhất cần đến lời khai đếm: An và Chi đều vô tội, nên
``hàng 1 có đúng một tội phạm'' ép Bình là tội phạm, và truy vấn
$\mathrm{SAT}(\mathit{KB}_2 \land \neg x_1)$ bị lan truyền đơn vị bác bỏ ngay
mà không cần phân nhánh. Toàn bộ $45$ lần gọi bộ giải trên bản đồ này đều được
giải chỉ bằng lan truyền đơn vị, khớp với cột phân nhánh bằng $0$ trong
Bảng~\ref{tab:benchmark_raw}.
\end{sloppypar}
